%=======================================================================
\section{Feature-Extraction Strategy}

No single descriptor captures everything that distinguishes the three tumor types. The pipeline
therefore uses a \textbf{multi-branch (multi-descriptor) design}: four complementary feature families
are computed independently and then fused. The intuition is that their errors are partly
uncorrelated, so the union is more discriminative than any one alone.

\begin{center}
\renewcommand{\arraystretch}{1.25}
\begin{tabularx}{\textwidth}{@{}l l Y Y@{}}
\toprule
\thd{Branch} & \thd{Descriptor} & \thd{Captures} & \thd{Question it answers} \\
\midrule
\textbf{A --- Stats}   & First-order statistics       & Global intensity distribution   & How bright / heterogeneous is the tissue overall? \\
\textbf{B --- Texture} & GLCM + LBP                   & Spatial texture \& micro-patterns & How are intensities arranged relative to each other? \\
\textbf{C --- DWT}     & Wavelet sub-band stats       & Multi-scale frequency content   & What structure appears at coarse vs.\ fine scales? \\
\textbf{D --- HOG}     & Gradient orientation hist.   & Edge \& shape structure         & What are the dominant edge directions and contours? \\
\bottomrule
\end{tabularx}
\renewcommand{\arraystretch}{1.0}
\end{center}

Each branch is extracted with its \textbf{own optimal hyper-parameters} (Section~\ref{sec:hyper}),
concatenated with first-order statistics, and the resulting long vector is cleaned and standardised
before modelling. Because the branches are produced separately, the pipeline can also benchmark them
individually and in combinations (A, B, C, D, A+B, B+C, A+B+C, and the full fusion) to show exactly
how much each contributes.

\begin{notebox}[Design principle:]
Diversity over redundancy. GLCM/LBP describe texture, DWT describes scale/frequency, HOG describes
shape/edges, and Stats describe global intensity --- four different physical views of the same image.
\end{notebox}
